% -*- mode: latex -*-
%
% Notations communes à la fiche d'exercices et à ses éléments de correction.
%
% Les définitions reprennent à l'identique celles de cours/preamble.tex
% (l. 110-123), de façon que les documents de travaux dirigés emploient
% exactement les mêmes symboles que les transparents, sans avoir à charger tout
% le préambule beamer (qui suppose \coursfile défini et appelle git via
% \write18).
%
% Convention structurante du cours, à respecter dans tous les énoncés : le
% « modèle de la nature » (le DGP) s'écrit avec des lettres grecques (\beta,
% \varepsilon), le « modèle de l'économètre » avec des latines grasses
% (\mathbf b) et des erreurs \epsilon. L'estimateur des MCO est \hat{\mathbf b}.
%
% Il n'existe pas de macro \E ni \Var dans ce dépôt : l'espérance et la variance
% s'écrivent \mathbb E[...] et \mathbb V[...]. Ne pas en introduire.

\newcommand{\trace}{\mathrm{tr}}
\newcommand{\tracarg}[1]{\mathrm{tr}\left\{#1\right\}}
\newcommand{\iid}[2]{\mathrm{iid}\left(#1,#2\right)}
\newcommand{\normal}[2]{\mathcal N\left(#1,#2\right)}
\newcommand{\epsvar}{\sigma_{\varepsilon}^2}
\newcommand{\sample}{\mathcal Y_T}
\newcommand{\samplet}[1]{\mathcal Y_{#1}}

% Convergences. Le cours indexe \plim et \dlim sur T ; les exercices écrits en
% coupe individuelle ont besoin de la variante en N, d'où les deux jeux.
\newcommand{\plimT}{\xrightarrow[T\rightarrow\infty]{\text{proba}}}
\newcommand{\dlimT}{\xrightarrow[T\rightarrow\infty]{\text{loi}}}
\newcommand{\plimN}{\xrightarrow[N\rightarrow\infty]{\text{proba}}}
\newcommand{\dlimN}{\xrightarrow[N\rightarrow\infty]{\text{loi}}}
